% Part: incompleteness
% Chapter: incompleteness-provability
% Section: lob-thm

\documentclass[../../../include/open-logic-section]{subfiles}

\begin{document}

\olfileid{inc}{inp}{tar}

\olsection{सत्याची अपरिभाष्यता}

\emph{परिभाष्यता} ही संकल्पना
अंकगणिताच्या भाषेसाठी
आकारिक चिन्हार्थमीमांसा
उपलब्ध असण्यावर
अवलंबून असते.
अंकगणिताच्या भाषेतील
सूत्रे आणि वाक्ये
यांचा एक संच आपण
वर्णिला आहे.
ही वाक्ये नैसर्गिक
संख्यांबद्दल विधाने
करतात असे वाचणे
हे त्यांचे “अभिप्रेत
अर्थनिर्धारण” आहे;
असे विधान सत्य
किंवा असत्य असू
शकते.
$\Struct{N}$ ही
$\Nat$ आधारक्षेत्र
असलेली आणि
अंकगणिताच्या भाषेतील
चिन्हांचे मानक
अर्थनिर्धारण करणारी
!!{structure} असू द्या.
मग $\Sat{N}{!A}$
याचा अर्थ “$!A$
मानक अर्थनिर्धारणात
सत्य आहे” असा होतो.

\begin{defn}
नैसर्गिक संख्यांवरील
$R(x_1,\dots,x_k)$ हा
संबंध $\Struct{N}$
मध्ये \emph{परिभाष्य}
असतो, जर आणि फक्त जर
अंकगणिताच्या भाषेत
असे $!A(x_1,\dots,x_k)$
सूत्र असते की
प्रत्येक $n_1,\dots,n_k$
साठी $R(n_1,\dots,n_k)$
हे $\Sat{N}{!A(\num n_1,\dots,\num
  n_k)}$ असेल
तर आणि तरच खरे असते.
\end{defn}

दुसऱ्या शब्दांत,
संबंध $\Struct{N}$
मध्ये परिभाष्य
असतो, जर आणि फक्त
जर तो $\Th{TA}$
या उपपत्तीत निरूपणीय
असतो; येथे
$\Th{TA} =
\Setabs{!A}{\Sat{N}{!A}}$
हा अंकगणितातील
सत्य वाक्यांचा
संच आहे.
(हे लगेच स्पष्ट
झाले नसेल, तर
व्याख्या पुन्हा
तपासून याची
खात्री करून घ्या.)

\begin{lem}
प्रत्येक संगणनक्षम
संबंध $\Struct{N}$
मध्ये परिभाष्य
असतो.
\end{lem}

\begin{proof}
$\Th{Q}$ मध्ये
एखाद्या संबंधाचे
निरूपण करणारे
सूत्र $\Struct{N}$
मध्ये त्याच संबंधाची
परिभाषा देते, हे
सहज तपासता येते.
\end{proof}

आता उलट विधानही
खरे आहे का, असे
विचारता येते:
$\Struct{N}$ मध्ये
परिभाष्य असलेला
प्रत्येक संबंध
संगणनक्षम असतो
का? उत्तर नाही.
उदाहरणार्थ:

\begin{lem}
थांबण्याचा संबंध
$\Struct{N}$ मध्ये
परिभाष्य असतो.
\end{lem}

\begin{proof}
क्लिनीचे प्रमाण
रूप प्रमेय सांगते
की प्रत्येक आंशिक
संगणनक्षम फलन~$f$
साठी असा निर्देशांक~$e$
असतो की सर्व
$x \in \Nat$ साठी
$f(x) =
U(\umin{s}{T(e,x,s)})$;
येथे $U$ आणि $T$
आदिम पुनरावर्ती
आणि म्हणून सर्वत्र
परिभाषित आहेत.
त्यामुळे $f(x)$
परिभाषित असते
(म्हणजे संगणन
थांबते), जर आणि
फक्त जर असा~$s$
असतो की $T(e,x,s)$
खरे असते.

आता $H$ हा
थांबण्याचा संबंध
असू द्या, म्हणजे
\[
H = \Setabs{\tuple{e,x}}{\lexists[s][T(e, x, s)]}.
\]
$!D_T$ हे
$\Struct{N}$ मध्ये
$T$ ची परिभाषा
देऊ द्या. मग
\[
H = \Setabs{\tuple{e,x}}{\Sat{N}{\lexists[s][!D_T(\num e, \num x, s)]}},
\]
म्हणून $\lexists[s][!D_T(z, x, s)]$
हे $\Struct{N}$
मध्ये~$H$ ची
परिभाषा देते.
\end{proof}

\begin{prob}
$Q(n) \defiff n \in \Setabs{\Gn{!A}}{\Th{Q} \Proves !A}$
हा संबंध अंकगणितात
परिभाष्य आहे
हे दाखवा.
\end{prob}

मग $\Th{TA}$
चे काय? ती
अंकगणितात परिभाष्य
आहे का? म्हणजे,
$\Setabs{\Gn{!A}}{\Sat{N}{!A}}$
हा संच अंकगणितात
परिभाष्य आहे
का? तार्स्कीचे
प्रमेय याचे
नकारार्थी उत्तर
देते.

\begin{thm}
\ollabel{thm:tarski}
अंकगणितातील सत्य
!!{sentence}s चा संच
अंकगणितात परिभाष्य
नाही.
\end{thm}

\begin{proof}
$!D(x)$ त्याची
परिभाषा देते असे
समजू; म्हणजे
$\Sat{N}{!A}$ हे
$\Sat{N}{!D(\gn{!A})}$
असेल तर आणि
तरच खरे असते.
स्थिरबिंदू उपपत्तीवरून
असे बंद सूत्र $!A$
असते की
$\Th{Q} \Proves !A \liff \lnot !D(\gn{!A})$,
आणि म्हणून
$\Sat{N}{!A \liff \lnot !D(\gn{!A})}$.
पण मग
$\Sat{N}{!A}$ हे
$\Sat{N}{\lnot !D(\gn{!A})}$
असेल तर आणि
तरच खरे असते;
$!D(y)$ अंकगणितातील
सत्य विधानांचा
संच परिभाषित करते
या गृहीतकाला यामुळे
विरोध होतो.
\end{proof}

तार्स्कीने हे
विश्लेषण सत्याच्या
अधिक व्यापक
तत्त्वज्ञानात्मक
संकल्पनेला लावले.
कोणतीही भाषा $L$
दिल्यास, तार्स्कीच्या
मते $L$ साठी
सत्याची पर्याप्त
संकल्पना तिच्या
प्रत्येक वाक्य
$X$ साठी पुढील
अट पूर्ण करायला
हवी:
\begin{quote}
‘$X$’ सत्य आहे
जर आणि फक्त जर
$X$.
\end{quote}
इंग्रजीसाठी
तार्स्कीने दिलेले
वारंवार उद्धृत
केले जाणारे
उदाहरण मराठीत
असे मांडता येते:
\begin{quote}
‘बर्फ पांढरा आहे’
हे सत्य आहे
जर आणि फक्त जर
बर्फ पांढरा आहे.
\end{quote}
परंतु विकर्ण
फलनाचे निरूपण
करण्याइतकी प्रबळ
असलेली कोणतीही
भाषा आणि कोणतेही
भाषिक विधेय $T(x)$
घेतल्यास, “$X$
जर आणि फक्त जर
$T(\text{‘$X$’})$
नाही” असे
वाक्य $X$
बांधता येते.
सत्यविषयक विधेयाने
एकाच वाक्याला
सत्य आणि असत्य
दोन्ही ठरवावे,
असे आपल्याला
नको आहे.
म्हणून भाषेच्या
मर्यादेबाहेर
काही ना काही
प्रकारे पाऊल
टाकल्याशिवाय
त्या भाषेतील
सर्व वाक्यांसाठी
सत्यविषयक विधेय
निर्दिष्ट करता
येत नाही,
असा तार्स्कीने
निष्कर्ष काढला.
दुसऱ्या शब्दांत,
अशा पुरेशा प्रबळ भाषेतील
सर्व वाक्यांचे सत्यविषयक विधेय
त्याच भाषेत
परिभाषित करता
येत नाही.

\end{document}
